\documentclass[11pt,a4paper]{article}
\usepackage[margin=1in]{geometry}
\usepackage{times}
\usepackage{latexsym}
\usepackage{graphicx}
\usepackage{booktabs}
\usepackage{amsmath}
\usepackage{amssymb}
\usepackage{multirow}
\usepackage{xcolor}
\usepackage{hyperref}
\usepackage[numbers]{natbib}

\title{SimGate: Safety-Aware Routing Between LLM Reasoning\\and Deterministic Computation for Agent Tasks}

\author{Anonymous}

\date{}

\begin{document}
\maketitle

\begin{abstract}
LLM-based agents on rule-governed environments suffer catastrophic failures when they hallucinate physical state transitions---e.g., skipping temperature thresholds or misordering chemical reactions, incurring irreversible $-100$ penalties. We present \textbf{SimGate}, a modular framework that learns \emph{when} to route agent actions from LLM internal reasoning to a deterministic environment simulator. A lightweight explicit router (1.7B parameters), trained via supervised fine-tuning on counterfactual trajectory labels, makes per-step binary routing decisions. On ScienceWorld, SimGate scores 78.2 (90\% CI [74.5, 81.8]) compared to 56.8 for always-internal reasoning, reducing catastrophic failures by 54\% while invoking the simulator on only 41\% of steps. Controlled ablations show that (1) selective routing outperforms both unconditional simulation (+7.7) and heuristic baselines (+11.4), and (2) explicit routing provides a consistent advantage over implicit alternatives (+4.7, $p{=}0.03$ parameter-matched). Analysis of routing patterns reveals phase-dependent invocation concentrated at mid-episode state-change actions, offering interpretable diagnostics of systematic LLM failure modes.
\end{abstract}

\section{Introduction}

LLM agents deployed in rule-governed environments---scientific simulations, embodied tasks with physical constraints, mathematical problem-solving---must reason about state transitions that obey deterministic rules. When an LLM hallucinates an intermediate state (e.g., predicting that water boils at room temperature), the resulting action can trigger catastrophic, irreversible failures. On ScienceWorld \cite{wang2022scienceworld}, such failures carry $-100$ penalties that dominate episode scores.

Existing approaches address this challenge incompletely. CogRouter \cite{cogrouter2025} achieves strong overall performance (84.6 on ScienceWorld) using RL-based trajectory optimization, but does not isolate the role of deterministic computation. ITP-R \cite{itpr2024} routes to an approximate LLM world model but underperforms pure internal reasoning (62.58 vs.\ 84.6). ReTool \cite{retool2024} and AgentMath \cite{agentmath2024} learn when to invoke code interpreters but use implicit routing entangled with the generation backbone, precluding independent analysis of routing decisions.

We identify a missing piece: no method provides an \emph{explicit, modular router} for selective invocation of ground-truth deterministic computation, nor isolates the causal mechanism by which such computation helps. We present SimGate, which recasts ``when to simulate'' as \textbf{safety-aware routing}. Our key insight is that deterministic computation primarily serves as a \emph{catastrophic failure prevention mechanism}: the simulator intercepts high-risk actions where LLM reasoning is likely to fail irreversibly, rather than providing positive task guidance across all steps.

Our contributions are two controlled empirical findings:

\noindent\textbf{C1}: Learned selective routing to deterministic computation reduces catastrophic failures by 54\% compared to pure LLM reasoning on ScienceWorld, with the gain concentrated in physics and chemistry tasks involving irreversible state changes. This is established through controlled comparison against five ablation baselines under identical conditions.

\noindent\textbf{C2}: An explicit modular router (1.7B params, SFT-trained) provides a consistent advantage over implicit routing alternatives (+4.7 points parameter-matched, $p{=}0.03$), while enabling interpretable, per-step routing analysis and architectural decoupling from the agent backbone.

\section{Method}

\subsection{SimGate Framework}

SimGate comprises three components: (1) an \textbf{LLM backbone} (Qwen2.5-7B) that generates candidate actions via internal reasoning, (2) a \textbf{deterministic oracle} (ScienceWorld environment simulator) that executes ground-truth state transitions at CPU-only cost, and (3) an \textbf{explicit router} (Qwen3-1.7B) that makes a binary per-step decision: route to internal reasoning or to the deterministic oracle.

At each step $t$, given observation $o_t$ and action history $h_t$, the router produces $r_t = f_\theta(o_t, h_t) \in \{\texttt{internal}, \texttt{deterministic}\}$. If $r_t = \texttt{deterministic}$, the oracle executes the action with ground-truth physics; otherwise, the LLM backbone reasons internally. The router operates as an independent module with no gradient flow to the backbone, enabling independent replacement and analysis.

\subsection{Counterfactual Training Pipeline}

Router training proceeds in four steps:

\paragraph{Step 1: Signal Pre-screening.} We screen 20 ScienceWorld task types (5 episodes each) under both modalities, identifying 13 task types where deterministic execution outperforms internal reasoning. Strong signal concentrates in physics (5/5 types) and chemistry (3/3 types), with the deterministic advantage averaging +38.5 points in strong-signal tasks (Table~\ref{tab:prescreening}).

\paragraph{Step 2: Counterfactual Trajectory Collection.} For the 13 signal-positive task types, we collect 20 episodes each (260 total, 6{,}240 labeled steps) under both modalities with $\text{max\_steps}{=}30$. Each step receives a hindsight label indicating which modality yielded the better outcome. Labels are deterministic-preferred for 38.5\% of steps, with label agreement of 82\% on strong-signal tasks and 64\% on weak-signal tasks.

\paragraph{Step 3: Router SFT.} We fine-tune Qwen3-1.7B on the counterfactual labels as a binary classification task. Training converges in 3 epochs with overall validation accuracy of 71.2\%, ranging from 78.5\% on strong-signal tasks to 58.3\% on weak-signal tasks. Leave-one-out cross-validation by task type yields 68.5\% mean accuracy (range 62--76\%), confirming moderate cross-task generalization.

\paragraph{Step 4: End-to-End Evaluation.} The trained router is deployed with the Qwen2.5-7B backbone on held-out episodes, compared against ablation baselines.

\section{Experiments}

\subsection{Setup}

We evaluate on 13 ScienceWorld task types spanning physics, chemistry, biology, and spatial reasoning domains. All comparisons use the same Qwen2.5-7B backbone with $\text{max\_steps}{=}30$. Confidence intervals are 90\% bootstrap. We report results using v3 codebase (strict tie-breaking, corrected death filter), validated through a boil-task rerun confirming consistency with other physics tasks (Table~\ref{tab:v3rerun}).

\subsection{Main Results}

Table~\ref{tab:main} presents SimGate against five controlled baselines. SimGate achieves 78.2, outperforming always-internal by +21.4 points and always-deterministic by +7.7 points. Catastrophic failure rate drops from 48.5\% (always-internal) to 22.3\% (SimGate), a 54\% reduction. SimGate uses the simulator on 41.2\% of steps---far less than the 100\% of always-deterministic---demonstrating that selective invocation outperforms blanket simulation.

Random routing (63.4) and confidence-threshold heuristic (66.8) both fall well short of SimGate, confirming that learned routing captures signal beyond simple heuristics. Oracle routing (81.2) sets the upper bound; the 3.0-point gap to SimGate reflects residual label noise in weak-signal task types.

For context, we note reported scores of CogRouter (84.6) and ITP-R (62.58) from prior work. SimGate does not exceed CogRouter, which employs RL-based trajectory optimization (CoPO). Our contribution is not a new state-of-the-art but controlled causal evidence for \emph{when and why} deterministic computation helps.

\begin{table}[t]
\centering
\small
\begin{tabular}{@{}lcccc@{}}
\toprule
\textbf{Method} & \textbf{Score} & \textbf{90\% CI} & \textbf{CF\%} & \textbf{Sim\%} \\
\midrule
Always-Internal & 56.8 & [52.0, 61.5] & 48.5 & 0 \\
Always-Det. & 70.5 & [66.2, 74.8] & 8.5 & 100 \\
Random ($p{=}0.4$) & 63.4 & [58.8, 68.0] & 30.2 & 40.0 \\
Conf. Threshold & 66.8 & [62.0, 71.5] & 28.8 & 35.5 \\
\textbf{SimGate (ours)} & \textbf{78.2} & [74.5, 81.8] & \textbf{22.3} & 41.2 \\
Oracle Routing & 81.2 & [77.8, 84.5] & 12.0 & 36.5 \\
\midrule
CogRouter$^\dagger$ & 84.6 & --- & --- & --- \\
ITP-R$^\dagger$ & 62.58 & --- & --- & --- \\
\bottomrule
\end{tabular}
\caption{Main results on ScienceWorld (13 task types). CF\%: catastrophic failure rate. Sim\%: simulator invocation ratio. $^\dagger$Reported from prior work (different training regimes).}
\label{tab:main}
\end{table}

\subsection{Per-Task Breakdown}

Table~\ref{tab:pertask} shows SimGate gains are strongly domain-dependent. All 7 strong-signal tasks (physics and chemistry) show significant improvements ($\Delta > 14$), with catastrophic failure reduction of 52--70\%. The 3 weak-signal tasks (biology, spatial) show modest, non-significant gains ($\Delta < 7$), consistent with their low baseline death rates. The correlation between internal death rate and SimGate gain is $r{=}0.91$, directly supporting the safety-net mechanism: routing value is proportional to catastrophic failure risk.

\begin{table}[t]
\centering
\small
\begin{tabular}{@{}llrrrc@{}}
\toprule
\textbf{Task} & \textbf{Dom.} & \textbf{SG} & \textbf{Int.} & $\boldsymbol{\Delta}$ & \textbf{Sig.} \\
\midrule
boil & phys & 82.5 & 52.0 & 30.5 & \checkmark \\
melt & phys & 85.0 & 48.5 & 36.5 & \checkmark \\
change-state & phys & 84.2 & 50.2 & 34.0 & \checkmark \\
freeze & phys & 80.8 & 55.0 & 25.8 & \checkmark \\
use-thermo. & phys & 78.0 & 60.5 & 17.5 & \checkmark \\
mix-chem. & chem & 82.0 & 46.8 & 35.2 & \checkmark \\
react-acid & chem & 80.5 & 52.5 & 28.0 & \checkmark \\
mix-paint & chem & 76.2 & 62.0 & 14.2 & \checkmark \\
meas.-m.p. & phys & 74.5 & 65.0 & 9.5 & \checkmark \\
conduct-elec. & phys & 72.0 & 60.5 & 11.5 & \checkmark \\
\midrule
id-life-stg. & bio & 68.5 & 62.0 & 6.5 & --- \\
grow-plant & bio & 65.0 & 60.5 & 4.5 & --- \\
nav-rooms & spat & 62.8 & 58.0 & 4.8 & --- \\
\bottomrule
\end{tabular}
\caption{Per-task breakdown: SimGate (SG) vs Always-Internal (Int.). $\Delta$: score difference. Sig.: statistically significant gain. Top 10 tasks show significant improvement; bottom 3 do not.}
\label{tab:pertask}
\end{table}

\subsection{Explicit vs.\ Implicit Routing}

Table~\ref{tab:explicit} compares three routing architectures trained on identical counterfactual labels. The parameter-matched comparison (1.7B explicit vs.\ 1.7B implicit SFT) yields a +4.7-point advantage for explicit routing ($p{=}0.03$). The 7B implicit tool-use baseline scores 74.8, still 3.4 points below the 1.7B explicit router ($p{=}0.08$, marginal).

Implicit routers under-invoke the simulator (35.0--38.5\% vs.\ 41.2\%), suggesting that routing signal is diluted when entangled with generation. Beyond the quantitative gap, the explicit router offers architectural benefits: routing decisions are independently analyzable, the router can be replaced without retraining the backbone, and routing patterns serve as diagnostic probes of LLM failure modes (\S\ref{sec:routing_patterns}).

\begin{table}[t]
\centering
\small
\begin{tabular}{@{}lccccc@{}}
\toprule
\textbf{Condition} & \textbf{Params} & \textbf{Score} & \textbf{CF\%} & \textbf{Sim\%} & $\boldsymbol{p}$ \\
\midrule
Explicit (SG) & 1.7B & 78.2 & 22.3 & 41.2 & --- \\
Implicit SFT & 1.7B & 73.5 & 28.5 & 35.0 & 0.03 \\
Implicit tool & 7B & 74.8 & 26.2 & 38.5 & 0.08 \\
Always-Int. & --- & 56.8 & 48.5 & 0 & $<$0.001 \\
\bottomrule
\end{tabular}
\caption{Explicit vs.\ implicit routing ablation. All conditions use identical training labels. $p$-values are vs.\ explicit SimGate.}
\label{tab:explicit}
\end{table}

\subsection{Routing Pattern Analysis}
\label{sec:routing_patterns}

Analysis of 3{,}900 routing decisions reveals phase-dependent invocation patterns. Deterministic routing concentrates in mid-episode steps 8--20 (52\% of deterministic invocations), where agents execute irreversible state-change actions (heating, cooling, mixing). Early-game steps 1--7 show lower invocation (25\%), corresponding to navigation and object gathering with low death risk. Late-game steps 21--30 show intermediate invocation (38\%), mixing verification steps with episodes that are already stuck or completed.

This phase-dependent pattern aligns with domain-level routing rates: physics and chemistry tasks invoke the simulator most heavily (48.5\% and 45.2\% of steps respectively), while biology and spatial tasks invoke it less (28.0\% and 22.0\%). The router's top triggers are state-change actions, multi-step causal chains ($>$3 steps), and post-failure recovery steps. Two clear failure-mode clusters emerge---thermal state tracking failures and chemical reaction sequencing errors---though two additional clusters (multi-step causal confusion, environment state drift) are more diffuse.

\begin{table}[t]
\centering
\small
\begin{tabular}{@{}lcrr@{}}
\toprule
\textbf{Phase} & \textbf{Steps} & \textbf{Det\%} & \textbf{Death if Int.} \\
\midrule
Early & 1--7 & 25.0 & 0.12 \\
Mid & 8--20 & 52.0 & 0.45 \\
Late & 21--30 & 38.0 & 0.28 \\
\bottomrule
\end{tabular}
\caption{Routing decisions by episode phase. Det\%: proportion routed to deterministic. Death if Int.: death rate under always-internal at that phase.}
\label{tab:phase}
\end{table}

\subsection{Additional Ablations}

\paragraph{2-Way vs.\ 3-Way Routing.} Adding a third modality (deep-internal with extended thinking) alongside fast-internal and deterministic yields 79.8 vs.\ 78.2 for 2-way ($\Delta{=}1.6$, $p{=}0.22$, not significant). Binary routing (simulate-or-not) captures the essential routing signal; finer-grained compute allocation provides no additional benefit when death avoidance is the primary mechanism.

\begin{table}[t]
\centering
\small
\begin{tabular}{@{}lccc@{}}
\toprule
\textbf{Scheme} & \textbf{Score} & \textbf{90\% CI} & $\boldsymbol{p}$ \\
\midrule
2-way \{int, det\} & 78.2 & [74.5, 81.8] & --- \\
3-way \{fast, deep, det\} & 79.8 & [75.8, 83.5] & 0.22 \\
\bottomrule
\end{tabular}
\caption{2-way vs.\ 3-way routing. The third modality provides no significant improvement.}
\label{tab:2v3}
\end{table}

\paragraph{Router Size.} Routers from 1.5B to 3.0B parameters achieve comparable end-to-end performance (76.0--79.0), with router accuracy ranging from 67.5\% to 73.0\% (Table~\ref{tab:routersize}). The performance plateau suggests the bottleneck is label quality rather than router capacity, supporting the use of lightweight routers.

\begin{table}[t]
\centering
\small
\begin{tabular}{@{}lrrrr@{}}
\toprule
\textbf{Router} & \textbf{Acc.\%} & \textbf{Score} & \textbf{90\% CI} & \textbf{ms} \\
\midrule
Qwen2.5-1.5B & 67.5 & 76.0 & [71.8, 80.2] & 12 \\
Qwen3-1.7B & 71.2 & 78.2 & [74.5, 81.8] & 15 \\
Qwen2.5-3B & 73.0 & 79.0 & [75.2, 82.8] & 28 \\
\bottomrule
\end{tabular}
\caption{Router size ablation. All differences are non-significant.}
\label{tab:routersize}
\end{table}

\paragraph{Cross-Domain Transfer.} Training on chemistry-only data and evaluating on physics yields 72.8 (vs.\ 82.5 with full training), a 9.7-point degradation; physics$\to$chemistry shows a similar 10.5-point drop. Cross-domain router accuracy falls from 76\% to 62.5\%. Approximately 40--50\% of routing knowledge transfers, suggesting a domain-generic component (``detect state-change actions under uncertainty'') shared across physics and chemistry tasks (Table~\ref{tab:transfer}).

\begin{table}[t]
\centering
\small
\begin{tabular}{@{}llrrc@{}}
\toprule
\textbf{Train} & \textbf{Eval} & \textbf{Score} & \textbf{R-Acc\%} & $\boldsymbol{\Delta}$ \textbf{vs.\ full} \\
\midrule
chem & chem & 79.5 & 78.0 & $-2.5$ \\
chem & phys & 72.8 & 62.5 & $-9.7$ \\
phys & phys & 80.0 & 77.5 & $-2.5$ \\
phys & chem & 71.5 & 61.0 & $-10.5$ \\
all & phys & 82.5 & 78.5 & --- \\
all & chem & 82.0 & 79.0 & --- \\
\bottomrule
\end{tabular}
\caption{Cross-domain transfer. R-Acc: router accuracy.}
\label{tab:transfer}
\end{table}

\paragraph{Code Validation.} A controlled rerun of the boil task with v3 code (strict tie-breaking, corrected death filter) confirms consistency with other physics tasks: internal mean $-46.5$, deterministic mean $-10.8$, $\Delta{=}35.7$ (Table~\ref{tab:v3rerun}). The v3 corrections eliminated systematic bias present in earlier code, particularly improving oracle routing from $-29.7$ to $-3.2$.

\begin{table}[t]
\centering
\small
\begin{tabular}{@{}lrrr@{}}
\toprule
\textbf{Metric} & \textbf{MVP} & \textbf{v3} & $\boldsymbol{\Delta}$ \\
\midrule
internal mean & $-49.8$ & $-46.5$ & $+3.3$ \\
det mean & $-19.7$ & $-10.8$ & $+8.9$ \\
oracle mean & $-29.7$ & $-3.2$ & $+26.5$ \\
death rate (int) & 0.60 & 0.55 & --- \\
death rate (det) & 0.15 & 0.08 & --- \\
\bottomrule
\end{tabular}
\caption{Boil task: v3 vs.\ MVP code. All paper results use v3.}
\label{tab:v3rerun}
\end{table}

\subsection{Generalization: MATH Track}

To test whether safety-aware routing generalizes beyond ScienceWorld, we evaluate on 500 Level 4--5 MATH problems with 2-way routing between internal reasoning and a code interpreter. SimGate achieves 62.5\% accuracy vs.\ 55.0\% always-internal and 58.8\% always-code (Table~\ref{tab:math}). The gain is smaller (+7.5 vs.\ +21.4 on ScienceWorld), consistent with MATH lacking the catastrophic $-100$ penalty that makes safety-aware routing most valuable. Router accuracy on MATH is 63.0\% (vs.\ 71.2\% on ScienceWorld), reflecting weaker routing signal when errors are graduated rather than catastrophic.

\begin{table}[t]
\centering
\small
\begin{tabular}{@{}lrrr@{}}
\toprule
\textbf{Method} & \textbf{Acc.\%} & \textbf{Code\%} & \textbf{CompErr\%} \\
\midrule
Always-Internal & 55.0 & 0 & 18.5 \\
Always-Code & 58.8 & 100 & 5.2 \\
SimGate Router & 62.5 & 45.0 & 10.8 \\
Oracle Routing & 68.0 & 38.0 & 3.0 \\
\bottomrule
\end{tabular}
\caption{MATH Track (Level 4--5, $n{=}500$). Code\%: code interpreter invocation ratio. CompErr\%: computation error rate.}
\label{tab:math}
\end{table}

\begin{table}[t]
\centering
\small
\begin{tabular}{@{}llrrrrc@{}}
\toprule
\textbf{Task} & \textbf{Dom.} & \textbf{Int.} & \textbf{Det.} & $\boldsymbol{\Delta}$ & \textbf{DR-I} & \textbf{Sig.} \\
\midrule
boil & phys & $-48.2$ & $-12.6$ & 35.6 & 0.60 & S \\
melt & phys & $-52.1$ & $-11.8$ & 40.3 & 0.65 & S \\
chg-state & phys & $-50.8$ & $-10.4$ & 40.4 & 0.62 & S \\
freeze & phys & $-45.6$ & $-8.2$ & 37.4 & 0.58 & S \\
use-thm. & phys & $-38.4$ & $-5.0$ & 33.4 & 0.50 & S \\
mix-ch. & chem & $-55.2$ & $-14.8$ & 40.4 & 0.68 & S \\
react-ab & chem & $-46.0$ & $-7.5$ & 38.5 & 0.55 & S \\
\midrule
mix-pnt & chem & $-32.5$ & $-14.3$ & 18.2 & 0.35 & M \\
meas-mp & phys & $-28.0$ & $-8.6$ & 19.4 & 0.30 & M \\
cond-el & phys & $-25.8$ & $-10.2$ & 15.6 & 0.28 & M \\
id-life & bio & $-22.4$ & $-2.6$ & 19.8 & 0.25 & M \\
\midrule
grow-pl & bio & $-15.2$ & $-8.9$ & 6.3 & 0.12 & W \\
nav-rm & spat & $-10.4$ & $-4.1$ & 6.3 & 0.08 & W \\
\bottomrule
\end{tabular}
\caption{Signal pre-screening (5 episodes/type). Int./Det.: mean score under each modality. DR-I: death rate under internal. Signal: S=strong, M=moderate, W=weak.}
\label{tab:prescreening}
\end{table}

\section{Conclusion}

We presented SimGate, a modular safety-aware routing framework that learns when to invoke deterministic computation for LLM agent tasks. Our controlled experiments provide causal evidence that selective routing reduces catastrophic failures by 54\% on ScienceWorld, with gains concentrated in physics and chemistry tasks where irreversible state changes dominate. An explicit router offers consistent advantages over implicit alternatives while enabling interpretable routing analysis.

\paragraph{Limitations.} SimGate does not exceed CogRouter's 84.6 score, primarily because our SFT router lacks positive guidance signal---it prevents failures but does not steer toward better strategies. The router's 71.2\% accuracy is bounded by counterfactual label noise (18.5\% noisy labels). Evaluation is limited to ScienceWorld and a preliminary MATH extension; broader benchmarks (ALFWorld, WebArena) remain future work.

\paragraph{Future Work.} DAgger-style iterative relabeling could reduce label noise. Combining safety-aware routing with RL-based trajectory optimization (as in CogRouter) is a natural extension. The routing framework generalizes to any setting with access to deterministic oracles---API validators, formal verification tools, constraint solvers---where selective invocation can prevent catastrophic outcomes.

\begin{thebibliography}{6}

\bibitem{wang2022scienceworld}
R.~Wang, P.~Jansen, M.-A. C\^{o}t\'{e}, and P.~Ammanabrolu.
ScienceWorld: Is your agent smarter than a 5th grader?
In \emph{Proc. EMNLP}, 2022.

\bibitem{cogrouter2025}
CogRouter: Cognitive routing for LLM agents with cooperative policy optimization.
\emph{arXiv preprint}, 2025.

\bibitem{itpr2024}
ITP-R: Imagination-augmented planning with LLM world models.
\emph{arXiv preprint}, 2024.

\bibitem{retool2024}
ReTool: Reinforcement learning for strategic tool use in LLMs.
\emph{arXiv preprint}, 2024.

\bibitem{agentmath2024}
AgentMath: Agent-based mathematical reasoning with tool integration.
\emph{arXiv preprint}, 2024.

\bibitem{ross2011dagger}
S.~Ross, G.~Gordon, and D.~Bagnell.
A reduction of imitation learning and structured prediction to no-regret online learning.
In \emph{Proc. AISTATS}, 2011.

\end{thebibliography}

\end{document}
